\documentclass[11pt,a4paper,toc=bibliography,toc=listof,titlepage=firstiscover]{scrreprt}
\usepackage[utf8]{inputenc}
\usepackage{array}
\usepackage{blindtext}
\usepackage{graphicx}
\usepackage{subcaption}
\usepackage[T1]{fontenc}
\usepackage{setspace}
\usepackage{pdfsync}
\usepackage{blindtext}
\usepackage[ngerman]{babel}
\usepackage{booktabs}
\usepackage{xcolor}
\usepackage{multirow}
%Schriftarten hier definieren
\usepackage{lmodern}
%\usepackage{tgheros}
%\renewcommand*\familydefault{\sfdefault}
%\usepackage{tgpagella}
\usepackage[babel,german=guillemets]{csquotes}
\usepackage[style=apa]{biblatex}
\usepackage{hyperref}
\usepackage{csquotes}
\addbibresource{resources/literatur.bib}
\hypersetup{colorlinks=true, linkcolor=black, urlcolor=blue, citecolor=blue}
%Abstand der Einträge im Literaturverzeichnis als em
\setlength\bibitemsep{0.75em}
%Kein Einzug bei einem neuen Absatz
\setlength{\parindent}{0mm}
%Inhaltsverzeichnis ohne eigene Seitenzahl
\AtBeginDocument{\addtocontents{toc}{\protect\thispagestyle{empty}}} 
\usepackage[left=3cm,right=2.5cm,top=2.5cm,bottom=2cm]{geometry}

%Daten des Deckblatts
  \titlehead
  {
    {
      \begin{center}
        \vspace{2.5cm}
        \includegraphics[scale=0.3]{resources/nak-logo.png}
      \end{center}
    }
  	%M.Sc. Angewandte Informatik
  }
  \subject{Hausarbeit}
  \title{Einflussfaktoren auf die Schulleistungen im Rahmen der Oberschule}
    \subtitle{MADS2100 Reporting und Visualisierung 23oB}
    \author{Leon Henne}
  \date{\small{Köln, den \today}}
  \publishers{Betreut durch Dr. Robert Stahlbock}
\begin{document}
%Titelseite erzeugen
\maketitle
%Der Nachfolgende Text wird 1,5-zeilig gesetzt
\onehalfspacing
%Römische Ziffern für Abbildungsverzeichnis
\pagenumbering{Roman}
\setcounter{page}{0}
%Inhaltsverzeichnis erzeugen
\tableofcontents
\clearpage
%Abbildungs und Tabellenverzeichnis
\listoftables
\listoffigures
%ToDo --> Abkürzungsverzeichnis
\clearpage
\pagenumbering{arabic}
\setcounter{page}{1}
%Datei aus dem Ordner /chapter/ einbinden. So kann man mehrere Textdateien zu einem großen Dokument zusammenfügen.
\setcounter{chapter}{0}

\chapter[Problemstellung]{Problemstellung}

Der für jede Nation erstrebenswerte langfristige ökonomische Fortschritt wird unter anderem stark durch das vorherrschende Bildungsniveau beeinflusst (\cite[S. 1]{Cortez2008UsingDM}). 
Um dieses sich ergebende Bildungsniveau stärker zu durchleuchten und final zu verbessern, wird zur Unterstützung der Schüler und Lehrkräfte die Modellierung von Schulleistungen eingesetzt. \cite[S. 1]{Cortez2008UsingDM} 
So kann eine zeitabhängige Vorhersage der Leistungen, lernschwächere Schüler detektieren und damit Lehrkräfte frühzeitig befähigen, mit entsprechenden Maßnahmen in den Lernprozess einzugreifen. \cite[S. 2]{Namoun.2021} 
Verstärkt wurde dieser Bedarf durch die in der Vergangenheit eingetretenen Covid-Pandemie, und den damit verbundenen Schulschließungen, welche für neue erhebliche Herausforderungen sorgten. \cite[S. 2]{Clark.2021} 
Durch \cite[S. 13]{Clark.2021} konnte hierzu aufgezeigt werden, welchen positiven Effekt digitale Lernunterstützungen auf die Schülergruppen ausmachten.
Aus der Arbeit von \cite[S. 9]{Namoun.2021} geht jedoch hervor, dass bereits seit 2017 erneut das Interesse anstieg hinsichtlich der Modellierung von Lernergebnissen.
Seitdem besteht besonders ein Fokus auf die Untersuchung des Bildungsniveaus von Bachelorstudiengängen, sodass die Untersuchung weiterführender Schulen lediglich ein Anteil von in etwa 12\% besitzt. \cite[S. 11]{Namoun.2021} 
Die in den letzten Jahren erforschte Modellierung von Studierendenergebnissen lässt häufig unbeachtet, wie einzelne Faktoren innerhalb der maschinellen Lernmethoden zu den Vorhersagen führen. \cite[S. 19]{Namoun.2021}
Die Gesamtheit dieser aktuellen Gegebenheiten motiviert die nachfolgende Untersuchung des gewählten Datensatzes anhand der daran abgeleiteten Forschungsfragen.
\chapter[Zielsetzung]{Zielsetzung}

Inhalt dieses Kapitels ist die zum einen die Erläuterung des gewählten Datensatzes.
Zum Anderen werden daraus abgeleitete Forschungsfragen formuliert, welche im Rahmen dieser Arbeit durch die Entwicklung und Analyse von Visualisierungen untersucht werden. 

\section[Untersuchter Datensatz]{Datensatz}

\section[Forschungsfragen]{Forschungsfragen}
\chapter[Grundlagen]{Grundlagen}

\section[how to Visualisierungen ?]{how to Visualisierungen ?}
\section[how to analyse Visualisierungen ?]{how to analyse Visualisierungen ?}

\chapter[Fazit]{Fazit}


\begin{itemize}
    \item Lernleistung lediglich durch Noten bemessen
    \item alter Datensatz
    \item Komplexität ist ok, da durch das digitale format Visualisierungen in den angezeigten Daten gefiltert werden können. (zoom, kategorie auswahl)
\end{itemize}

%Literaturverzeichnis einfügen
\renewcommand*{\UrlFont}{\rmfamily}
\printbibliography
\section*{Eidesstattliche Erklärung}

Hiermit erkläre ich, dass ich die vorliegende Arbeit eigenständig und ohne fremde Hilfe angefertigt habe. Textpassagen, die wörtlich oder dem Sinn nach auf Publikationen oder Vorträgen anderer Autoren beruhen, sind als solche kenntlich gemacht.

Die Arbeit wurde bisher keiner anderen Prüfungsbehörde vorgelegt und auch noch nicht veröffentlicht. \newline

Köln, den \today

\vspace{15mm}

Leon Henne


\end{document}
